
\pagestyle{empty}
\newgeometry{margin=1cm}
\footnotesize

\begin{landscape}
\part*{Hannah Fry - Hello World}
\section{Themenübersicht}
Untenan finden Sie eine Übersicht einiger (Unter-)Themen, die im Buch von Hannah Fry behandelt werden. Es handelt sich hierbei um eine nicht abschliessende Auflistung. Weitere Themen können in Absprache mit der Lehrperson vereinbart werden. Beachten Sie zudem, dass die Themen nicht zwingend Unter-Kapiteln im Buch entsprechen, sondern dass gewisse Themen über ein Kapitel ``verteilt'' immer wieder auftauchen.
\label{sec:themen}

% Please add the following required packages to your document preamble:
% \usepackage{booktabs}
% \usepackage{multirow}
%\begin{table}[H]
	\begin{longtable}{@{}l|p{5.3cm}||p{10cm}|p{10cm}@{}}
		\toprule
		\textbf{Kapitel}                       & \textbf{Fragen / Thema}                        & \textbf{Auftrag}                                                                                                                                                                                                                                                                                                                                                                 & \textbf{Recherche-Auftrag                                                                                                                                                                                                                                                      }\\ \midrule\midrule
		\multirow{3}{*}{Macht}        & Algorithmen: Definition \& Kategorien & \begin{tabular}[c]{@{}l@{}}Was ist ein Algorithmus? Wie wird er definiert?\\ Welche vier Kategorien von Algorithmen nennt das Buch?\\ Präsentieren Sie die Kategorien, die im Buch erwähnt werden\end{tabular}                                                                                                                                                          & Finden Sie für jede im Buch erwähnte Algorithmus-Kategorie ein weiteres, eigenes Beispiel                                                                                                                                                                              \\\cmidrule(l){2-4}
		& Ethische Folgen von Algorithmen       & Welche negativen Folgen werden im Zusammenhang mit dem Umgang mit Algorithmen genannt? In welchen Lebensbereichen / Situationen sollten Algorithmen nicht oder mit Vorsicht angewandt werden?                                                                                                                                 & Führen Sie eine eigene Recherche zu heiklen Algorithmen durch, präsentieren Sie Ihre Resultate der Klasse.                                                                                                                                                             \\\cmidrule(l){2-4} 
		& Eigene Einschätzung \& Fazit          & Fassen Sie die Kernaussagen des Kapitels zusammen. Was nehmen Sie vom Lesen des Kapitels mit? Welche anderen Aspekte sind Ihnen beim Lesen aufgefallen / geblieben? Was scheint Ihnen wichtig? & Recherchieren Sie, inwiefern ChatGPT einem der erwähnten Algorithmen zugeordnet werden kann oder nicht. Haben wir es mit einer neuen Kategorie von Algorithmus zu tun?                                                                                                 \\\midrule\midrule
		\multirow{5}{*}{Daten}        & Wert von Daten                        & Welche Aspekte von Daten werden im Kapitel als wertvoll bezeichnet?                                                                                                                                                                                                                                                                                                     & Recherchieren Sie die 5 wertvollsten Firmen der Welt und Beschreiben Sie, inwiefern Daten für diese Firmen eine Rolle spielen.                                                                                                                                         \\\cmidrule(l){2-4} 
		& Datenschutz                           & Was sind anonymisierte Daten? Welche Informationen können gemäss Buch selbst aus anonymisierten Daten herausgelesen werden?                                                                                                                                                                                                                                             & Recherchieren Sie, was das Internet über Sie weiss. Suchen Sie insbesondere nach ``Data Breaches'' in der Schweiz. Präsentieren Sie die Ergebnisse der Klasse                                                                                                            \\\cmidrule(l){2-4} 
		& Mikromanipulation \& Psychologie       & Was ist Mikro-Manipulation? Wie werden wir im Alltag durch Daten (mikro-)manipuliert?                                                                                                                                                                                                                                                                                   & Finden Sie 5 Lebensbereiche, in denen Sie manipuliert werden.                                                                                                                                                                                                          \\\cmidrule(l){2-4} 
		& Politische Einflussnahme              & Was sagt das Buch dazu, inwiefern Daten politische Prozesse beeinflussen können?                                                                                                                                                                                                                                                                                        & Recherchieren und Präsentieren Sie die Rolle von Cambridge Analytica im Trump-Wahlkampf.                                                                                                                                                                               \\\cmidrule(l){2-4} 
		& Bewertungen \& Bürger-Rating          & Fassen Sie Ihre Erkenntnisse sowie die Aussagen aus dem Buch zum Thema Bewertungen zusammen Wie werden Menschen durch Firmen und Regierungen bewertet? Welche Interessen verfolgen diese dabei?                                                                                                                                                                                                                                            & Recherchieren Sie, in welchen weiteren Formen Menschen freiwillig oder unfreiwillig durch Firmen oder Regierungen bewertet werden (z.B. FICO-Score, Fichen-Skandal etc.).\\\midrule\midrule
		\multirow{4}{*}{Justiz}       & Algorithmen in der Justiz             & Wie werden Entscheidungsbäume und \textit{Random Forests} durch das Buch beschrieben? Fassen Sie die Beispiele aus dem Buch zu diesem Thema zusammen.                                                                                                                                                                                  & Recherchieren Sie, was ein Entscheidungsbaum und was ein Random Forest ist. Fassen Sie zusammen, wie diese funktionieren und wie diese angewandt werden. Finden Sie ein weiteres Anwendungsbeispiel ausserhalb der Justiz.                                             \\\cmidrule(l){2-4} 
		& Justiz \& Datenschutz                 & Fassen Sie die Datenschutz-Regelungen zusammen, welche im Buch erwähnt werden. Was ist an juristischen Daten anders / sensibler als an anderen Daten? Inwiefern betreffen Sie diese Regelungen?                                                                                                                                                                         & Welche neuen Datenschutz-Regeln sind seit dem Erscheinen des Buches dazugekommen? Führen Sie Ihre eigene Recherche durch!                                                                                                                                              \\\cmidrule(l){2-4} 
		& Falsche und richtige Urteile          & Stellen Sie die Problematik der ``falschen-positiven'', ``richtig-positiven'', ``falsch-negativen'' und ``richtig-negativen'' Vorhersagen / Urteile vor und beurteilen Sie das Thema in Bezug auf die Justiz. Erklären Sie ebenfalls die Begriffe ``Spezifizität'' und ``Sensitivität''. Stellen Sie die Problematik der ``maschinellen Vorurteile'' klar.                            & Recherchieren Sie mindestens 2 Beispiele von falsch-positiven oder falsch-negativen Entscheidungen durch die AI. Beschreiben Sie, ob diese Fälle generelle Zweifel an der AI aufkommen lassen, oder ob es sich um Einzelfälle handelt.                                 \\\cmidrule(l){2-4} 
		& Schwierige Entscheide                 & Illustrieren Sie das Problem der ``schwierigen Entscheide'' anhand der Denkensweisen, die im Buch ``Schnelles Denken, langsames Denken'' von Daniel Kahneman erwähnt werden (dieses wird im Buch von Fry kurz angesprochen, Sie müssen es nicht lesen).                                                                                                                     & Finden Sie weitere Beispiele von schwierigen Aufgaben, die ``kontra-intuitives'' Denken zur Findung der Lösung erfordern. Mit ``kontra-intuitivem Denken'' ist langsames, analytisches Denken gemeint. Bonus: können solche Aufgaben durch heutige AI gut bewältigt werden? \\\midrule\midrule
		\multirow{3}{*}{Medizin}      & Mustererkennung                       & Erklären Sie, wie maschinelle Mustererkennung in der heutigen Medizin angewandt wird, indem Siedie Lern- und Funktionsweise von ``Neuronalen Netzwerken'' erklären.                                                                                                                                                                                                       & Illustrieren Sie, wofür Konvolutive Neuronale Netzwerke (Convolutional Neural Networks, CNNs) auch ausserhalb der Medizin genutzt werden können                                                                                                                        \\\cmidrule(l){2-4} 
		& Krebsdiagnosen \& AI                  & Erklären Sie, wie KI zur besseren Entdeckung von Krebszellen beitragen kann und welche Herausforderungen sich betreffend Über- und Unter-Diagnose stellen. Erklären Sie ebenfalls die Begriffe ``sensitiv'' und ``sepezifisch'' (s. Kapitel zu Justiz). Welche Herausforderungen stellen sich beim Trainieren von KI für medizinische Diagnosen betreffend Trainings-Daten? & Finden Sie beispiele von Schweizer Spitälern, die KI bereits heute zu Krebsdiagnosezwecken verwenden.                                                                                                                                                                  \\\cmidrule(l){2-4} 
		& Medizin \& Datenschutz                & Fassen Sie die Datenschutz-Problematiken zusammen, welche im Buch erwähnt werden. Was ist an medizinischen Daten anders / sensibler als an anderen Daten? Inwiefern könnten Sie diese Regelungen persönlich betreffen?  Erklären Sie die Problematik anhand kostenpflichtiger Gentests (Herkunfts-Tests).                                                               & Finden Sie ein Beispiel aus der Schweiz, in welchem die digitale Speicherung von medizinischen Personendaten zu Bedenken aus Sicht des Datenschutz führt.                                                                                                              \\\midrule\midrule
		\multirow{2}{*}{Autos}        & Autos \& Algorithmen                  & Erklären Sie, wie der Algorithmus von Bayes funktioniert und wie dieser für selbstfahrende Autos verwendet wird. Erklären Sie ebenfalls, weshalb es so schwierig ist, einen Algorithmus für selbstfahrende Autos ``manuell'' zu programmieren.                                                                                                                           & Finden Sie weitere Beispiele, in welchen Bayes angewandt wird. Recherchieren Sie zudem, ob Bayes auch heute noch angewandt wird für selbstfahrende Autos.                                                                                                              \\\cmidrule(l){2-4} 
		& Autos \& Ethik                        & Erläutern Sie einige ethische Problematiken, die mit selbstfahrenden Autos verbunden sind. Erklären Sie insbesondere die Problematik schwieriger ethischer Entscheidungen im Falle eines sicheren, nahenden Unfalls (FahrerIn oder PassantIn schützen?) sowie die Problematik abnehmender Fahrer-Fähigkeiten.                                                           & Finden Sie mindestens 5 weitere Problematiken, die mit selbstfahrenden Autos verbunden sind. Sofern Sie Ihre Überlegungen schlüssig und überzeugend darlegen, dürfen Sie hier auch Ihre eigenen Überlegungen anstellen statt zu recherchieren.                         \\\midrule\midrule
		\multirow{2}{*}{Kriminalität} & Kriminalität \& Algorithmen           & Erläutern Sie, wie Kriminalität mittels PredPol bekämpft wird und beschreiben Sie die Funktionsweise von PredPol. Welche Vorgänger-Systeme gab es, inwiefern haben diese versagt, und weshalb sind gerade räumliche (geografische) Daten besonders aufschlussreich?                                                                                                     & Recherchieren Sie, welche Polizei-IT-Systeme in Europa und in der Schweiz angewandt werden, um nach Kriminellen zu fahnden.                                                                                                                                            \\\cmidrule(l){2-4} 
		& Kriminalität \& Falschdiagnosen       & Erläutern Sie detailliert die Problematik der doubles und fassen Sie die Gesichtserkennungs-Algorithmen zusammen, welche im Buch erwähnt werden.                                                                                                                                                                                                                        & Recherchieren Sie, wie sich der Social Credit Score / Sesame Score in China seit Veröffentlichung des Buches weiterentwickelt hat und präsentieren Sie die Resultate ausführlich (min. 3-4 Minuten).                                                                   \\\midrule\midrule
		\multirow{2}{*}{Kunst}        & Fake Kunst \& Manipulation            & Erläutern Sie das Thema ``\textit{social proof}'' anhand der im Buch aufgeführten Beispiele für Streaming-Dienste.                                                                                                                                                                                                                                                                 & Finden Sie 2 weitere Beispiele für social proof mit Bezug zur Informatik. Dabei kann es sich auch um Beispiele aus der Finanzwelt handeln (z.B. Bitcoin \& Social Proof, Game Stop Hype \& Social Proof etc.).                                                         \\\cmidrule(l){2-4} 
		& Kunst \& Qualität                     & Erläutern Sie die Problematik der Messung von Qualität in der Kunst anhand der Beispiele aus dem Buch.                                                                                                                                                                                                                                                                  & Präsentieren Sie mindestens 3 weitere Beispiele aus der Kunst und erstellen Sie, wenn möglich, eigene Beispiele (z.B. Suno, Meta Imagine etc.). Sind die Beispiele originell / von hoher Qualität?                                                                     \\ 
	\end{longtable}
	
%\end{table}
\end{landscape}


\restoregeometry
\pagestyle{fancy}
\newpage
\section{Auftrag}
Lesen Sie im Buch ``Hello World'' von Hannah Fry die folgenden Kapitel: 
\begin{enumerate*}
	\item ``Einleitung''
	\item ``Macht''
	\item ``Daten''
	\item Ein Kapitel zu einem Thema  Ihrer Wahl (s. Abschnitt \ref{sec:themen}).
\end{enumerate*}
Selbstverständlich dürfen Sie auch weitere Kapitel aus dem Buch lesen. Machen Sie sich Notizen zu jedem Kapitel (worum geht es? Was ist Ihnen geblieben? Was verstehen Sie / nicht? etc.). 

Im Anschluss an die Lektüre werden Sie folgende Abgaben in einer 3er- bis 4er-Gruppe erstellen:

\begin{enumerate}
	\item Ein \textbf{Fact Sheet} zum Kapitel Ihrer Wahl
	\item Eine \textbf{Präsentation} Ihres Kapitels.
\end{enumerate}

Das \textit{Fact Sheet} (Punkt 1), sowie die Präsentation (Punkt 2) zählen beide je zu 50\% und machen gemeinsam \textit{eine schriftliche Note} aus. Die Anforderungen für jede dieser Abgaben sind in den folgenden Abschnitten kurz erklärt.

\subsection{Fact Sheet}
Bilden Sie Dreier- oder Vierergruppen nach Kapitel, das Sie präsentieren. Pro Gruppe bearbeiten Sie also genau eines der sieben Kapitel. Jedes Kapitel wird von genau einer Gruppe bearbeitet. 

\subsubsection*{Austausch in der Gruppe}
Tauschen Sie sich in der Gruppe zu folgenden Themen aus:

\begin{itemize}
	\item \textbf{Gedankenaustausch zum Gelesenen}: Worum ging es? Was ist Ihnen \textit{persönlich} geblieben?
	\item \textbf{Organisation des Auftrags}: evtl. Fragen zum Vortrag, Ideen-Sammlung, Vorgehen, Aufgabenverteilung 
\end{itemize}


Erstellen Sie nun in der Gruppe zu Ihrem Thema ein Fact Sheet. Es besteht aus 2 A4-Seiten und enthält die relevanten Informationen zu Ihrem Thema, auf einfache und verständliche Weise zusammengefasst. Als Ergänzung zu Ihren kurzen Zusammenfassungs-Texten sollten Sie \textit{sinnvolle} Bilder, Grafiken und/oder Diagramme erstellen. Das Fact Sheet enthält zudem kontroverse Fragestellungen, welche die Grundlage für die Gruppendiskussionen (s. Abschnitt \ref{subsec:ind}) bilden.

\textbf{Bewertungskriterien}:
\begin{itemize}
	\item (2 Pt) Inhaltlich korrekte und umfassende Zusammenfassung des Kapitels
	\item (1 Pt) Mindestens eine selbst erstellte \textbf{Illustration} (Bild, Diagramm, Grafik, usw.), die Ihr Factsheet sinnvoll ergänzt.
	\item (1 Pt) Formelle Kriterien
	\begin{itemize}
		\item Abgabeformat: PDF
		\item Dateiname: ``\texttt{[Klassenname] - K[Kapitelnummer] [Kapitelname] Factsheet.pdf}'' (z.B. ``\texttt{2d - K2 Daten Factsheet.pdf}'')
		\item 2 A4-Seiten
		\item Korrekte Sprache, vollständige Sätze
		\item Schriftgrösse ca. 11 Punkte
		\item Vor- und Nachnamen aller Gruppenmitglieder im PDF enthalten
	\end{itemize}
	\item (1 Pt) Formulieren Sie am Ende Ihrer Zusammenfassung zwei interessante, reichhaltige Diskussions-Fragen (die nicht einfach mit ja / nein beantwortet werden können) oder provokativ-pointierte Aussagen, die zur Diskussion anregen. 
\end{itemize}

Das Factsheet ist spätestens eine Woche vor dem Vortrag abzugeben.

\subsection{Vortrag}
\label{subsec:ind}

Erstellen Sie einen Vortrag in Gruppen zu Ihrem Kapitel  (s. Abschnitt \ref{sec:themen}):

\begin{itemize}
	\item Dauer der Präsentation: ca. \textbf{20 Minuten} (ohne Diskussion), die Redezeit der Gruppenmitglieder sollte dabei ungefähr gleich verteilt sein. Danach moderieren Sie eine kurze Diskussion (s. Abschnitt untenhalb dieser Liste)
	\item Der Inhalt Ihrer Präsentation soll eine didaktisch hochwertige Zusammenfassung des von Ihnen behandelten Kapitels aus dem Buch von Hannah Fry sein (s. Abschnitt \ref{sec:themen}). Die Unterthemen des Kapitels sollten alle behandelt werden, inklusive Recherche-Aufträge.
	\item Tragen Sie sich für ein Thema Ihrer Wahl ein (\href{https://eduzh-my.sharepoint.com/:x:/g/personal/cyril_wendl_edu_zh_ch/EScHfDarDeVHhAJB3WPCghsBYa63XpdaEkcrZk5dOiQiVA?e=kHy0bM}{Link}). Falls es zu viele InteressentInnen für ein Kapitel gibt, sprechen Sie sich bitte untereinander ab, oder kontaktieren Sie die Lehrperson, um eine mögliche Aufteilung des Kapitels in mehrere Unter-Kapitel zu besprechen.
	\item Die detaillierten Bewertungskriterien finden Sie auf \href{https://moodle.ksimlee.ch/mod/assign/view.php?id=32090}{Moodle}.
\end{itemize}

Im Anschluss an den Vortrag findet kurze Diskussionsrunde über die Fragen bzw. provokanten Aussagen in der Klasse statt, welche durch die vortragenden Personen moderiert wird.
